% =========================================================
% preambulo.tex
% Paquetes y configuración compartida por TODO el documento.
% Se carga una sola vez desde main.tex
% =========================================================



% ---------- Codificación e idiomas ----------
\usepackage[utf8]{inputenc}
\usepackage[T1]{fontenc}
% Toda la tesis se escribe en español (el cuerpo, el índice, "Capítulo",
% "Figura", "Tabla", etc.). Se deja "english" cargado por si en algún
% punto necesitas una cita o término puntual en inglés con \foreignlanguage.
% IMPORTANTE: el ÚLTIMO idioma listado es el que queda activo por defecto,
% por eso "spanish" va al final.
% La opción "es-tabla" fuerza que las tablas digan "Tabla X" en vez del
% default de babel, que es "Cuadro X".
\usepackage[english,spanish,es-tabla]{babel}
% ---------- Geometría de página (A4, como el PDF original) ----------
\usepackage[a4paper,margin=2.5cm]{geometry}

% ---------- Tipografía del CUERPO del documento ----------
% El PDF de Víctor usa "NimbusRomNo9L", el clon libre de Times
% que trae TeX Live. mathptmx activa esa fuente (texto y matemáticas).
\usepackage{mathptmx}

% ---------- Tipografía de la PORTADA (Gill Sans) ----------
% Se deja disponible aquí, pero SOLO se usa dentro de portada.tex
% con \sffamily, para no afectar el resto de los capítulos.
\usepackage{gillius2}

% ---------- Imágenes ----------
\usepackage{graphicx}
\usepackage{svg}
\graphicspath{{figuras/}} % así no repites la ruta en cada capítulo
\usepackage{subcaption}

% ---------- Color para tablas ----------
% La opción [table] habilita \rowcolor (paquete colortbl por dentro),
% para poder pintar filas completas de una tabla.
\usepackage[table]{xcolor}

% ------------------------------------------------
% Color institucional para tablas
% ------------------------------------------------
\definecolor{azultabla}{HTML}{4F81BC}

% Encabezado de tabla
\newcommand{\encabezadotabla}{%
	\rowcolor{azultabla}%
	\color{white}\bfseries
}

\newcommand{\encabezado}[1]{%
	\cellcolor{azultabla}\color{white}\bfseries #1%
}


% ---------- Matemáticas ----------
\usepackage{amsmath}
\usepackage{amssymb}

% ---------- Pseudocódigo / algoritmos ----------
% algorithm2e: fácil de traducir a español (SetKw, nombres de "Entrada",
% "Salida", "mientras", "si", etc.) y se ve limpio en tesis.
% [ruled,vlined]: caja con línea superior/inferior y líneas verticales
% de indentado, que es el estilo más común en tesis de cómputo.
\usepackage[ruled,vlined,linesnumbered]{algorithm2e}
\SetKwInput{KwEntrada}{Entrada}
\SetKwInput{KwSalida}{Salida}
\renewcommand{\algorithmcfname}{Algoritmo}

% ---------- Enlaces y referencias cruzadas ----------
\usepackage[numbers]{natbib}
\usepackage{hyperref}
\hypersetup{
    colorlinks=true,
    linkcolor=black,
    citecolor=black,
    urlcolor=blue,
    pdftitle={Genetic Non-Linear Polynomial Discriminant Analysis},
    pdfauthor={Ángel García Báez}
}

%---- Manejo de la bibliografia (BibTeX clásico numérico) ----%
% ---------- Bibliografía: BibTeX clásico + natbib + estilo DIN ----------
% natdin.bst necesita el paquete natbib para funcionar (no es opcional).
% [numbers] hace que las citas en el texto salgan como [1], [2]... en vez de (Autor, Año).
\bibliographystyle{natdin}
